% GENERATED FILE. Edit canonical lessons, then run npm run generate:books.
% canonical-source-hash: fnv1a64:69eefe06eab6aa0b
% canonical-lessons: FR-C325-historien, FR-C325-ancetre, FR-C325-patrimoine, FR-C325-archeologue, FR-C325-fouille

\chapter{Les traces du passé}
\label{ch:fr-les-traces-du-passe}
\hlchaptermodality{325}

\begin{chapteropening}
\textbf{By the end of this chapter:} \emph{I can talk about people who study the past and what they find: the historian, ancestors, heritage, the archaeologist and a dig.}

Complete the last lesson of chapter 325: I can talk about people who study the past and what they find: the historian, ancestors, heritage, the archaeologist and a dig.
\end{chapteropening}

\section[l'historien]{l'historien — the historian}

\label{lesson:FR-C325-historien}



\begin{quote}
Before the new one: say the French for the conquest, then the French for to defeat.
\end{quote}

\subsection*{You'll want to know: l'historien}
\textbf{l'historien} — \textquotedblleft{}the historian\textquotedblright{}.

\begin{quote}
Les historiens ne sont pas d'accord sur cette date.
\end{quote}

\begin{quote}
\emph{Historians do not agree on this date.}
\end{quote}

Feminine: \emph{l'historienne}.

\begin{cousinweb}
\emph{historien} is built on \emph{histoire}, from Latin \emph{historia}, a borrowing from Greek \emph{historía} \textquotedblleft{}inquiry, knowledge\textquotedblright{} (Lewis \& Short).
\end{cousinweb}

\begin{grammarlens}[title={What you've built}]
One more word for telling, explaining and describing at B1.
\end{grammarlens}

\subsection*{Your turn}
\begin{itemize}
\raggedright
  \item \emph{Say it:} \emph{l'historien}
  \item \emph{Say it:} \emph{l'historien}, once more
  \item \emph{Say it:} a sentence of your own with \emph{l'historien}
  \item \emph{Recall:} say the French for the conquest, then the French for to defeat, then say \emph{l'historien} again
\end{itemize}

\begin{tcolorbox}[breakable,colback=teal!4,colframe=teal!35!black,title={Before you move on}]
What does l'historien mean? (\textquotedblleft{}The historian\textquotedblright{}.) Say it once more.
\end{tcolorbox}
\section[l'ancêtre]{l'ancêtre — the ancestor}

\label{lesson:FR-C325-ancetre}



\begin{quote}
Before the new one: say the French for to defeat, then the French for the historian.
\end{quote}

\subsection*{You'll want to know: l'ancêtre}
\textbf{l'ancêtre} — \textquotedblleft{}the ancestor\textquotedblright{}.

\begin{quote}
Mes ancêtres vivaient dans ce village.
\end{quote}

\begin{quote}
\emph{My ancestors lived in this village.}
\end{quote}

Also a forerunner: \emph{l'ancêtre du vélo}.

\begin{cousinweb}
\emph{ancêtre} comes from Old French \emph{ancestre} (EtymDB), inherited from Latin \emph{antecessor} \textquotedblleft{}one who goes before\textquotedblright{}, from \emph{ante} and \emph{cēdere}.
\end{cousinweb}

\begin{grammarlens}[title={What you've built}]
One more word for telling, explaining and describing at B1.
\end{grammarlens}

\subsection*{Your turn}
\begin{itemize}
\raggedright
  \item \emph{Say it:} \emph{l'ancêtre}
  \item \emph{Say it:} \emph{l'ancêtre}, once more
  \item \emph{Say it:} a sentence of your own with \emph{l'ancêtre}
  \item \emph{Recall:} say the French for to defeat, then the French for the historian, then say \emph{l'ancêtre} again
\end{itemize}

\begin{tcolorbox}[breakable,colback=teal!4,colframe=teal!35!black,title={Before you move on}]
What does l'ancêtre mean? (\textquotedblleft{}The ancestor\textquotedblright{}.) Say it once more.
\end{tcolorbox}
\section[le patrimoine]{le patrimoine — heritage; an estate}

\label{lesson:FR-C325-patrimoine}



\begin{quote}
Before the new one: say the French for the historian, then the French for the ancestor.
\end{quote}

\subsection*{You'll want to know: le patrimoine}
\textbf{le patrimoine} — \textquotedblleft{}heritage; an estate\textquotedblright{}.

\begin{quote}
Ce château fait partie du patrimoine du pays.
\end{quote}

\begin{quote}
\emph{This castle is part of the country's heritage.}
\end{quote}

\emph{Le patrimoine culturel} — cultural heritage.

\begin{cousinweb}
\emph{patrimoine} is a learned borrowing from Latin \emph{patrimōnium} \textquotedblleft{}an inheritance from a father\textquotedblright{} (EtymDB), built on \emph{pater} \textquotedblleft{}father\textquotedblright{}.
\end{cousinweb}

\begin{grammarlens}[title={What you've built}]
One more word for telling, explaining and describing at B1.
\end{grammarlens}

\subsection*{Your turn}
\begin{itemize}
\raggedright
  \item \emph{Say it:} \emph{le patrimoine}
  \item \emph{Say it:} \emph{le patrimoine}, once more
  \item \emph{Say it:} a sentence of your own with \emph{le patrimoine}
  \item \emph{Recall:} say the French for the historian, then the French for the ancestor, then say \emph{le patrimoine} again
\end{itemize}

\begin{tcolorbox}[breakable,colback=teal!4,colframe=teal!35!black,title={Before you move on}]
What does le patrimoine mean? (\textquotedblleft{}Heritage; an estate\textquotedblright{}.) Say it once more.
\end{tcolorbox}
\section[l'archéologue]{l'archéologue — the archaeologist}

\label{lesson:FR-C325-archeologue}



\begin{quote}
Before the new one: say the French for the ancestor, then the French for heritage.
\end{quote}

\subsection*{You'll want to know: l'archéologue}
\textbf{l'archéologue} — \textquotedblleft{}the archaeologist\textquotedblright{}.

\begin{quote}
Les archéologues ont trouvé des pièces très anciennes.
\end{quote}

\begin{quote}
\emph{The archaeologists found some very old coins.}
\end{quote}

The science is \emph{l'archéologie}.

\begin{cousinweb}
\emph{archéologue} is a learned word built from Greek \emph{arkhaîos} \textquotedblleft{}ancient\textquotedblright{} and \emph{-logos} \textquotedblleft{}one who studies\textquotedblright{} (EtymDB lists no parent).
\end{cousinweb}

\begin{grammarlens}[title={What you've built}]
One more word for telling, explaining and describing at B1.
\end{grammarlens}

\subsection*{Your turn}
\begin{itemize}
\raggedright
  \item \emph{Say it:} \emph{l'archéologue}
  \item \emph{Say it:} \emph{l'archéologue}, once more
  \item \emph{Say it:} a sentence of your own with \emph{l'archéologue}
  \item \emph{Recall:} say the French for the ancestor, then the French for heritage, then say \emph{l'archéologue} again
\end{itemize}

\begin{tcolorbox}[breakable,colback=teal!4,colframe=teal!35!black,title={Before you move on}]
What does l'archéologue mean? (\textquotedblleft{}The archaeologist\textquotedblright{}.) Say it once more.
\end{tcolorbox}
\section[la fouille]{la fouille — the dig, the excavation; a search}

\label{lesson:FR-C325-fouille}



\begin{quote}
Before the new one: say the French for heritage, then the French for the archaeologist.
\end{quote}

\subsection*{You'll want to know: la fouille}
\textbf{la fouille} — \textquotedblleft{}the dig, the excavation; a search\textquotedblright{}.

\begin{quote}
Les fouilles ont duré trois ans.
\end{quote}

\begin{quote}
\emph{The dig went on for three years.}
\end{quote}

Also a search, by the police or at an airport.

\begin{cousinweb}
\emph{fouille} is built on \emph{fouiller} \textquotedblleft{}to dig, to search\textquotedblright{}, usually traced to Vulgar Latin \emph{fodiculāre} (EtymDB lists no parent), built on Latin \emph{fodere} \textquotedblleft{}to dig\textquotedblright{} (Lewis \& Short).
\end{cousinweb}

\begin{grammarlens}[title={What you've built}]
One more word for telling, explaining and describing at B1.
\end{grammarlens}

\subsection*{Your turn}
\begin{itemize}
\raggedright
  \item \emph{Say it:} \emph{la fouille}
  \item \emph{Say it:} \emph{la fouille}, once more
  \item \emph{Say it:} a sentence of your own with \emph{la fouille}
  \item \emph{Recall:} say the French for heritage, then the French for the archaeologist, then say \emph{la fouille} again
\end{itemize}

\begin{tcolorbox}[breakable,colback=teal!4,colframe=teal!35!black,title={Before you move on}]
What does la fouille mean? (\textquotedblleft{}The dig, the excavation; a search\textquotedblright{}.) Say it once more.
\end{tcolorbox}
